% Class-compatibility preamble retained for reader excerpts of
%   Hiraku Nakajima, "Quiver Varieties and Branching", SIGMA 5 (2009) 003,
%   DOI 10.3842/SIGMA.2009.003, CC BY 4.0.
%
% The source pins the publisher class sigma, which is not part of the delivery
% baseline, so this file supplies the amsart baseline and then copies the author's
% own declaration block mechanically and unchanged (pinned-original-source.tex
% lines 43-239): the aenume environment, the Theorem/Corollary/Lemma/Proposition and
% definition-style newtheorem declarations, the equation numbering, the author's
% reference macros (\lemref and the like) and the author's mathematical macros.
% Everything from pinned source line 240 onwards is left out: those lines define
% \rf, which uses the publisher class internal dimension \y@boxdim, and the two
% colour rules that serve that one macro; nothing in the delivered excerpts uses
% \rf, \graysquare or \halfgray.
% No mathematical macro, statement, symbol or assumption is introduced, changed or
% removed by this file.  The author's package mathdots (pinned source line 20) is not
% reproduced because it is not installed in the delivery TeX installation and no
% delivered excerpt uses its commands; the author's remaining packages are kept.
\documentclass[10pt,a4paper]{amsart}
\usepackage{amssymb}
\usepackage{amsbsy,amscd,mathtools}
\usepackage{mathrsfs}
\usepackage{nicefrac}
\usepackage[all]{xy}
\usepackage{color}
\usepackage{hyperref}
\newenvironment{aenume}{%
  \begin{enumerate}%
  \renewcommand{\theenumi}{\alph{enumi}}%
  \renewcommand{\labelenumi}{\textup{(\theenumi)}}%
  }{\end{enumerate}}
%

\usepackage{verbatim}
\usepackage{version}
\usepackage{color}
\newenvironment{NB}{
\color{red}{\bf NB}. \footnotesize
}{}
% For a final version, NB disappears.
\excludeversion{NB}
% Older version
%\renewenvironment{NB}{%
%  \comment
%  }{\endcomment}
% The following is used inside NB
\newenvironment{NB2}{
\color{blue}{\bf NB}. \footnotesize
}{}
%       Theorem environments
%% \theoremstyle{plain} %% This is the default
\newtheorem{Theorem}[equation]{Theorem}
\newtheorem{Corollary}[equation]{Corollary}
\newtheorem{Lemma}[equation]{Lemma}
\newtheorem{Proposition}[equation]{Proposition}
%\newtheorem{Axiom}{Axiom}

\theoremstyle{definition}
\newtheorem{Definition}[equation]{Definition}
\newtheorem{Example}[equation]{Example}
\newtheorem{CounterExample}[equation]{Counter-Example}
\newtheorem{Notation}[equation]{Notation}
\newtheorem{Convention}[equation]{Convention}
\newtheorem{Conjecture}[equation]{Conjecture}

%\theoremstyle{remark}
\newtheorem{Remark}[equation]{Remark}
\newtheorem{Remarks}[equation]{Remarks}
%\newtheorem*{Claim}{Claim}
%\renewcommand{\theClaim}{} % to make the claim environment unnumbered
\newtheorem{Question}[equation]{Question}
%\renewcommand{\theQuestion}{}  % to make the notation environment
                               % unnumbered
\newtheorem{Exercise}[equation]{Exercise}
%\renewcommand{\theExercise}{}

\numberwithin{equation}{section}
%\numberwithin{Theorem}{section}
%\numberwithin{figure}{section}

\newcommand{\thmref}[1]{Theorem~\ref{#1}}
\newcommand{\secref}[1]{Section~\ref{#1}}
\newcommand{\lemref}[1]{Lemma~\ref{#1}}
\newcommand{\propref}[1]{Proposition~\ref{#1}}
\newcommand{\corref}[1]{Corollary~\ref{#1}}
\newcommand{\subsecref}[1]{Section~\ref{#1}}
\newcommand{\chapref}[1]{Chapter~\ref{#1}}
\newcommand{\defref}[1]{Definition~\ref{#1}}
\newcommand{\remref}[1]{Remark~\ref{#1}}
\newcommand{\remsref}[1]{Remarks~\ref{#1}}
%
%       Math definitions
%\newcommand{\lsp}[2]{{}^{#1}{#2}}
\newcommand{\lsp}[2]{{\mskip-.3mu}{}^{#1}\mskip-1mu{#2}}
\newcommand{\defeq}{\overset{\operatorname{\scriptstyle def}}{=}}
\newcommand{\interior}{\rule{6pt}{.4pt}\rule{.4pt}{6pt}} % interior
                                                         % product
\newcommand{\C}{{\mathbb C}}
\newcommand{\Z}{{\mathbb Z}}
\newcommand{\Q}{{\mathbb Q}}
\newcommand{\R}{{\mathbb R}}
% Projective spaces
\newcommand{\proj}{{\mathbb P}}
%\newcommand{\CP}{\operatorname{\C P}}
\newcommand{\CP}{\proj}
\newcommand{\RP}{\operatorname{\R P}}
% Lie Groups
\newcommand{\SL}{\operatorname{\rm SL}}
\newcommand{\SU}{\operatorname{\rm SU}}
\newcommand{\GL}{\operatorname{GL}}
\newcommand{\PGL}{\operatorname{PGL}}
\newcommand{\U}{\operatorname{\rm U}}
\newcommand{\SO}{\operatorname{\rm SO}}
\newcommand{\grpSp}{\operatorname{\rm Sp}}
% Lie algebras
\newcommand{\algsl}{\operatorname{\mathfrak{sl}}} % because \sl="slant"
\newcommand{\su}{\operatorname{\mathfrak{su}}}
\newcommand{\gl}{\operatorname{\mathfrak{gl}}}
\newcommand{\algu}{\operatorname{\mathfrak{u}}}
\newcommand{\so}{\operatorname{\mathfrak{so}}}
\newcommand{\algsp}{\operatorname{\mathfrak{sp}}}
\newcommand{\g}{{\mathfrak g}}
\newcommand{\h}{{\mathfrak h}}
% Useful symbols
\newcommand{\Spec}{\operatorname{Spec}\nolimits}
\newcommand{\Proj}{\operatorname{Proj}\nolimits}
\newcommand{\End}{\operatorname{End}}
\newcommand{\Hom}{\operatorname{Hom}}
\newcommand{\Ext}{\operatorname{Ext}}
\newcommand{\Ker}{\operatorname{Ker}}
\newcommand{\Coker}{\operatorname{Coker}}
\newcommand{\Ima}{\operatorname{Im}}
\newcommand{\coker}{\operatorname{coker}}
\newcommand{\Endskew}{\operatorname{Endskew}}
\newcommand{\im}{\mathop{\text{\rm im}}\nolimits}
%\newcommand{\dim}{\mathop{\text{\rm dim}}\nolimits}
\newcommand{\rank}{\operatorname{rank}}
\newcommand{\ind}{\mathop{\text{\rm ind}}\nolimits}
\newcommand{\tr}{\operatorname{tr}}
\newcommand{\Tr}{\operatorname{Tr}}
\newcommand{\DB}{\overline{\partial}}
\newcommand{\pd}[2]{\frac{\partial#1}{\partial#2}}
\newcommand{\deriv}[1]{\frac{d}{d#1}}
\newcommand{\id}{\operatorname{id}}
\newcommand{\ve}{\varepsilon}
%
% notations used in this paper
%
% \newcommand{\vin}[1]{\operatorname{in}(#1)} % incoming vertex
% \newcommand{\vout}[1]{\operatorname{out}(#1)} % outgoing vertex
\newcommand{\vin}[1]{\operatorname{i}(#1)} % incoming vertex
\newcommand{\vout}[1]{\operatorname{o}(#1)} % outgoing vertex
\newcommand{\bM}{{\mathbf M}} % the vector space
\newcommand{\bN}{{\mathbf N}} % the vector space
\newcommand{\M}{{\mathfrak M}} % the moduli space
\newcommand{\Mreg}{\M^{\operatorname{s}}}
\newcommand{\La}{{\mathfrak L}} % the Lagrangian variety
\newcommand{\Fi}{{\mathfrak F}} % the fixed point variety
\newcommand{\dslash}{/\!\!/} % for algebro-geometric quotient (double slash)
\newcommand{\bv}{{\mathbf v}} % the vector v
\newcommand{\bw}{{\mathbf w}} % the vector w
\newcommand{\bC}{{\mathbf C}} % the Cartan matrix
\newcommand{\bA}{{\mathbf A}} % the adjacency matrix
\newcommand{\bI}{{\mathbf I}} % the identity matrix
\newcommand{\bn}{{\mathbf n}} % the kernel of the extended Cartan matrix
\newcommand{\bu}{{\mathbf u}}
\newcommand{\be}{{\mathbf e}}
\newcommand{\bq}{{\mathbf q}}
\newcommand{\bs}{\mathbf s}
\newcommand{\bc}{\mathbf c}
\newcommand{\codim}{\operatorname{codim}} % codimension
\newcommand{\gr}{\operatorname{gr}} % grade
\newcommand{\topdeg}{{\operatorname{top}}} % top degree
\newcommand{\Uq}{{\mathbf U}_q(\mathfrak g)} % the QUE algebra
\newcommand{\Un}{{\mathbf U}_q}
\newcommand{\Um}{\Tilde{\mathbf U}_q(\mathfrak g)} % the modified
                                                 % enveloping algebra
\newcommand{\Ua}{{\mathbf U}_q(\widehat{\mathfrak g})} % the quantum affine
                                                 % algebra
\newcommand{\Uae}{{\mathbf U}_{\varepsilon}(\widehat{\mathfrak g})} %
                                                 % the quantum affine
                                                 % algebra
\newcommand{\Ul}{{\mathbf U}_q({\mathbf L}{\mathfrak g})} % the quantum
                                                      % loop algebra
\newcommand{\Uma}{\Tilde{\mathbf U}^\bw_q(\widehat{\mathfrak g})}
%                                                 the quantum affine algebra
\newcommand{\Ui}{{\mathbf U}^{\Z}_q(\mathfrak g)}
\newcommand{\Ue}{{\mathbf U}_{\varepsilon}(\mathfrak g)}
\newcommand{\Ueone}{{\mathbf U}_{\varepsilon = 1}(\mathfrak g)}
\newcommand{\Uli}{{\mathbf U}^{\Z}_q({\mathbf L}\mathfrak g)}
\newcommand{\Ule}{{\mathbf U}_{\varepsilon}({\mathbf L}\mathfrak g)}
\newcommand{\bU}{\mathbf U} % the enveloping algebra
\newcommand{\Uqm}{{\widetilde{\mathbf U}_q}} % its modified form
\newcommand{\Pa}{{\mathfrak P}} % Hecke correspondence
\newcommand{\HomE}{\operatorname{E}}
\newcommand{\HomL}{\operatorname{L}}
\newcommand{\Zreg}{Z^{\operatorname{reg}\nolimits}}
\newcommand{\hiw}{{L}} % the highest weight module
%      \interval is used to provide better spacing after a [ that
%      is used as a closing delimiter.
\newcommand{\interval}[1]{\mathinner{#1}}
\newcommand{\shfO}{\mathcal O}
\newcommand{\Mp}{\widetilde\M}
\newcommand{\Mpp}{\widehat\M}
\newcommand{\Wedge}{{\textstyle \bigwedge}}
%
\newcommand{\Zw}{Z(\bw)}
\newcommand{\Mw}{\M(\bw)}
\newcommand{\Law}{\La(\bw)}
%\newcommand{\Zw}{Z(\bullet,\bullet;\bw)}
%\newcommand{\Mw}{\M(\bullet,\bw)}
%\newcommand{\Law}{\La(\bullet,\bw)}
\newcommand{\ngpiso}{\overset{.}{=}}
\newcommand{\idl}{\mathfrak I}
\newcommand{\Ktop}{K_{\operatorname{top}}}
\newcommand{\Kztop}{K_{0,\operatorname{top}}}
\newcommand{\Kotop}{K_{1,\operatorname{top}}}

\newcommand{\qH}{\mathbb H}
\newcommand{\D}{\mathcal D}
\newcommand{\A}{\mathcal A}
\newcommand{\incarrow}{\xrightarrow{\mathrm{inclusion}}}
\newcommand{\proarrow}{\xrightarrow{\mathrm{projection}}}
\usepackage{cleveref}
\crefname{originalequation}{original equation}{original equations}
\Crefname{originalequation}{Original equation}{Original equations}
\makeatletter
\newlabel{eq:bullet}{{5.2}{source}{}{sourcepointer-1}{}}
\newlabel{eq:bullet@cref}{{[originalequation][0][]5.2}{[1][0][]source}{}{}{}}
\makeatother
\crefname{originalequation}{original equation}{original equations}
\Crefname{originalequation}{Original equation}{Original equations}
\makeatletter
\newlabel{eq:branching}{{5.3}{source}{}{sourcepointer-2}{}}
\newlabel{eq:branching@cref}{{[originalequation][0][]5.3}{[1][0][]source}{}{}{}}
\makeatother
\crefname{originalequation}{original equation}{original equations}
\Crefname{originalequation}{Original equation}{Original equations}
\makeatletter
\newlabel{eq:Levistratum}{{2.17}{source}{}{sourcepointer-3}{}}
\newlabel{eq:Levistratum@cref}{{[originalequation][0][]2.17}{[1][0][]source}{}{}{}}
\makeatother
\crefname{originalequation}{original equation}{original equations}
\Crefname{originalequation}{Original equation}{Original equations}
\makeatletter
\newlabel{eq:basechange}{{2.18}{source}{}{sourcepointer-4}{}}
\newlabel{eq:basechange@cref}{{[originalequation][0][]2.18}{[1][0][]source}{}{}{}}
\makeatother
\title{The decomposition of the pushforward of the intersection cohomology sheaf of the regular part of a quiver variety}
\author{Hiraku Nakajima}
\date{doi:10.3842/SIGMA.2009.003; CC BY 4.0}
\begin{document}
\maketitle
\noindent\textit{Editorial note.} All delivered mathematical text is the author's own, copied from the pinned publisher source sources/algebra-fl01-quiver-varieties-and-branching/pinned-original-source.tex; no statement or proof step was rewritten, reordered or removed.\par\medskip
\noindent\textit{Editorial note.} Preamble provenance: this card uses sources/algebra-fl01-quiver-varieties-and-branching/retained-preamble.tex, which supplies the amsart baseline in place of the publisher class sigma and copies the author's declaration block of pinned source lines 43-239 unchanged, including the macros Hom, IC, Res and the theorem environments that the counted statement and proof use.\par\medskip
\noindent\textit{Editorial note.} Bibliography provenance: the counted proof cites section 10.2 of the author's reference Na-alg for the identification of the top-degree homology with an irreducible integrable representation on the regular strata. That citation is resolved with sources/algebra-fl01-quiver-varieties-and-branching/bibliography.tex, an excerpt of the author's own thebibliography environment (pinned source lines 4098-4257) copied character for character; no entry was rewritten and no citation key was changed.\par\medskip
\noindent\textit{Editorial note.} External references are navigation pointers only; every deliberately kept step of the author's argument stays inside the counted spans. Their printed numbers are the author's own numbering under the author's newtheorem and equation-numbering declarations, computed mechanically by the lane's scanner.\par\medskip
\noindent\textit{Editorial note.} Two displays inside the counted proof carry labels of the author (eq:nonregular and eq:restriction) that the counted text itself never cross-references, so they need no stub.\par\medskip
\noindent\textit{Allowed standard prerequisites (editorial).}\begin{itemize}
\item The local systems phi on the strata of M\_0(w) and the decomposition of the top-degree homology of the fibre as a direct sum of phi tensor V\_(G-hat),phi. Established in the previous subsection and used at the start of the counted proof (proof lines 2734-2738). Supplies the g-module decomposition whose summands the rest of the proof identifies.
\item The intersection cohomology sheaves IC of the strata and their behaviour under a finite birational morphism, in particular the projection formula used for kappa and kappa-prime. Used inside the counted proof at source lines 2766-2776. Justifies replacing the intersection cohomology of the closure of a product stratum by the pushforward of the exterior product under the finite birational morphism.
\item The geometric description of the regular strata: on the preimage of a regular stratum the map pi\_\{0,zeta\} is an isomorphism, and the top-degree homology of the fibre is the irreducible integrable highest weight representation with highest weight vector the class of x. Attributed by the author to section 10.2 of the cited paper Na-alg and used at proof lines 2742-2749. Identifies the two families of summands appearing in the final formula: the trivial-local-system case on the regular strata of M\_0 and the irreducible representation V\textasciicircum{}(g\_I\textasciicircum{}0 + h)(w - v\textasciicircum{}0) on the regular strata of M\_\{zeta\textasciicircum{}bullet\}.
\end{itemize}\medskip
\ifdefined\hypertarget\hypertarget{sourcepointer-1}{}\fi
\noindent\textit{Source reference.} Original equation 5.2: sources/algebra-fl01-quiver-varieties-and-branching/pinned-original-source.tex, source lines 2625--2625. Navigation pointer only: it labels the author's displayed decomposition of the pushforward that the counted proof refers to when it defines the multiplicities; the display itself is not reproduced.\par\medskip
\ifdefined\hypertarget\hypertarget{sourcepointer-2}{}\fi
\noindent\textit{Source reference.} Original equation 5.3: sources/algebra-fl01-quiver-varieties-and-branching/pinned-original-source.tex, source lines 2634--2634. Navigation pointer only: it labels a second displayed relation of the author that the proof substitutes into; the display is not reproduced.\par\medskip
\ifdefined\hypertarget\hypertarget{sourcepointer-3}{}\fi
\noindent\textit{Source reference.} Original equation 2.17: sources/algebra-fl01-quiver-varieties-and-branching/pinned-original-source.tex, source lines 1658--1658. Navigation pointer only: it labels the author's displayed product description of a stratum, which the counted proof recalls in one clause; the display is delivered verbatim inside the statement span of the card for prop:Levi but is not part of this card.\par\medskip
\ifdefined\hypertarget\hypertarget{sourcepointer-4}{}\fi
\noindent\textit{Source reference.} Original equation 2.18: sources/algebra-fl01-quiver-varieties-and-branching/pinned-original-source.tex, source lines 1738--1738. Navigation pointer only: it labels the author's commutative diagram used in the middle step; the diagram is not reproduced.\par\medskip
\begin{Theorem}\label{thm:branch}
  Let $V^\g(\lambda)$ \textup(resp.\ $V^{(\g_{I^0}+\mathfrak
    h)}(\lambda)$\textup) denote the irreducible integrable highest
  weight module of $\g$ \textup(resp.\ $\g_{I^0}+\mathfrak h$\textup)
  with the highest weight $\lambda$. Then we have
  \begin{gather*}
    (\pi_{0,\zeta^\bullet})_*(IC(\Mreg_{\zeta^\bullet}(\bv^0,\bw)))
   =
    \bigoplus_{\bv'}
    \Hom_{\g_{I^0}+\mathfrak
    h}(V^{(\g_{I^0}+\mathfrak h)}(\bw \!-\! \bv^0), V^\g(\bw\!-\!\bv'))
  \!\otimes\! IC(\Mreg_0(\bv',\bw))
  \\
  \phantom{(\pi_{0,\zeta^\bullet})_*(IC(\Mreg_{\zeta^\bullet}(\bv^0,\bw)))=}{} \oplus
  \bigoplus (\text{\rm
    $IC$ sheaves associated with {\it non\/}-regular strata}).
  \end{gather*}
\end{Theorem}

\begin{proof}
Let $\M_0(\bw)_{(\widehat G)}$ be a stratum of $\M_0(\bw)$. Take a point
$x\in \M_0(\bw)_{(\widehat G)}$. The decomposition
$H_{\topdeg}(\pi_{0,\zeta}^{-1}(x)) = \bigoplus
\phi\otimes V_{(\widehat G),\phi}$ in the previous subsection is a
decomposition of a $\g$-module.
%
Suppose $\M_0(\bw)_{(\widehat G)}$ is regular, i.e., $=
\Mreg_0(\bv',\bw)$ for some $\bv'$.
%
Then it is known that $H_{\topdeg}(\pi_{0,\zeta}^{-1}(x))$ is an
irreducible integrable representation of $\g$ with the highest weight
vector $[x] \in H_{\topdeg}(\pi_{0,\zeta}^{-1}(x)\cap
\M_\zeta(\bv',\bw))$ \cite[10.2]{Na-alg}, where
$x$ can be considered as a point in $\M_\zeta(\bv',\bw)$ as
$\pi_{0,\zeta}\colon \M_\zeta(\bv',\bw)\to \M_0(\bv',\bw)$ is an
isomorphism on the preimage of $\Mreg_0(\bv',\bw)$.
%
In particular, $V_{(\widehat G),\phi} = 0$ unless $\phi$ is the
trivial local system in this case.

Similarly if a stratum $\M_{\zeta^\bullet}(\bw)_{(\widehat H)}$ is of
the form $\Mreg_{\zeta^\bullet}(\bv^0,\bw)$ for some $\bv^0$, then
$H_{\topdeg}(\pi_{\zeta^\bullet,\zeta}^{-1}(y))$ is the irreducible
integrable representation $V^{(\g_{I^0}+\mathfrak h)}(\bw-\bv^0)$ if
$y\in \Mreg_{\zeta^\bullet}(\bv^0,\bw)$.
The highest weight vector is $[y]$ as above.

Consider a stratum $\M_{\zeta^\bullet}(\bw)_{(\widehat H)}$ which
contains $IC(\Mreg_0(\bv',\bw))$ in the decomposition
\eqref{eq:bullet}. We can write it as
$\Mreg_{\zeta^\bullet}(\bv^0,\bw)\times (\M_0)_{(\widehat H')}$ as in
\eqref{eq:Levistratum}. We consider the diagram
\eqref{eq:basechange}. Since $\kappa$ is a f\/inite birational morphism,
we have
\begin{equation*}
  IC(\overline{\Mreg_{\zeta^\bullet}(\bv^0,\bw)
     \times (\M_0)_{(\widehat H')}})
   = \kappa_*\big(
     IC(\M(\bv^0,\bw))\boxtimes IC(\overline{(\M_0)_{(\widehat H')}})
   \big),
\end{equation*}
and hence
\begin{equation*}\label{eq:nonregular}
  (\pi_{0,\zeta^\bullet})_* IC(\overline{\Mreg_{\zeta^\bullet}(\bv^0,\bw)
     \times (\M_0)_{(\widehat H')}}) = \kappa'_*\big(
   (\pi_{0,\zeta^\bullet})_*IC(\M(\bv^0,\bw))\boxtimes
   IC(\overline{(\M_0)_{(\widehat H')}})\big).
\end{equation*}
Since $\kappa'$ is also a f\/inite morphism, this contains only IC
sheaves on nonregular strata unless the second factor
$(\M_0)_{(\widehat H')}$ is trivial.
%
Hence \eqref{eq:branching} now becomes
\begin{equation}\label{eq:restriction}
   \operatorname{Res} V^{\g}(\bw - \bv')
   = \bigoplus_{\bv^0} V_{\bw-\bv'}^{\bw-\bv^0}
   \otimes V^{(\g_{I^0}+\mathfrak h)}(\bw - \bv^0),
\end{equation}
where $V_{\bw-\bv'}^{\bw-\bv^0}$ is the multiplicity of
$IC(\Mreg_0(\bv',\bw))$ in the decomposition of\linebreak[2]
$(\pi_{0,\zeta^\bullet})_*\linebreak[2](IC(\Mreg_{\zeta^\bullet}(\bv^0,
\linebreak[2]\bw)))$ in
\eqref{eq:bullet}. Thus
$V_{\bw-\bv'}^{\bw-\bv^0} =
\Hom_{\g_{I^0}+\mathfrak h}
(V^{(\g_{I^0}+\mathfrak h)}(\bw - \bv^0), V^\g(\bw-\bv'))$.
%
Note here that $V^{(\g_{I^0}+\mathfrak h)}(\bw - \bv^0)$ are
non-isomorphic for dif\/ferent $\bv^0$.
\end{proof}

% Bibliography excerpt retained for reader excerpts of
%   Hiraku Nakajima, "Quiver Varieties and Branching", SIGMA 5 (2009) 003,
%   DOI 10.3842/SIGMA.2009.003, CC BY 4.0.
%
% Every entry below is copied character for character from the author's own
% thebibliography environment of the pinned publisher source
% (pinned-original-source.tex lines 4098-4257).  No entry was rewritten,
% reordered, merged, split or invented, and no citation key was changed;
% only the entries that the delivered excerpts cite are kept.
\begin{thebibliography}{99}

\bibitem{CB}
Crawley-Boevey W., Geometry of the moment map for representations of
  quivers, {\it Compositio Math.} \textbf{126} (2001),   257--293.
 % \MR{MR1834739 (2002g:16021)}


\bibitem{Kac}
 Kac V.G., Inf\/inite-dimensional Lie algebras, 3rd ed., Cambridge
  University Press, Cambridge, 1990. %\MR{MR1104219 (92k:17038)}


\bibitem{KS}
Kashiwara M., Saito Y., Geometric construction of crystal bases, {\it Duke
  Math. J.} \textbf{89} (1997),  9--36,  \mbox{\href{http://arxiv.org/abs/q-alg/9606009}{q-alg/9606009}}. %\MR{MR1458969 (99e:17025)}


\bibitem{Na-quiver}
Nakajima H., Instantons on {ALE} spaces, quiver varieties, and
  Kac--Moody algebras, {\it Duke Math. J.} \textbf{76} (1994), 365--416.
  %\MR{MR1302318 (95i:53051)}


\bibitem{Na-proc}
Nakajima H., Varieties associated with quivers, in Representation Theory of
  Algebras and Related Topics (Mexico City, 1994), {\it CMS Conf. Proc.}, Vol.~19,
  Amer. Math. Soc., Providence, RI, 1996, 139--157.
  %\MR{MR1388562   (97m:16022)}


\bibitem{Na-alg}
Nakajima H., Quiver varieties and Kac--Moody algebras, {\it Duke Math. J.}
  \textbf{91} (1998), 515--560. %\MR{MR1604167 (99b:17033)}


\bibitem{Lecture}
Nakajima H., Lectures on Hilbert schemes of points on surfaces, {\it University
  Lecture Series}, Vol.~18, American Mathematical Society, Providence, RI, 1999.
  %\MR{MR1711344 (2001b:14007)}


\bibitem{Na-qaff}
Nakajima H., Quiver varieties and f\/inite-dimensional representations of
  quantum af\/f\/ine algebras, {\it J. Amer. Math. Soc.} \textbf{14} (2001),
  145--238, \href{http://arxiv.org/abs/math.QA/9912158}{math.QA/9912158}. %\MR{MR1808477 (2002i:17023)}


\bibitem{Na-tensor}
Nakajima H., Quiver varieties and tensor products, {\it Invent. Math.}
  \textbf{146} (2001), 399--449, \href{http://arxiv.org/abs/math.QA/0103008}{math.QA/0103008}. %\MR{MR1865400 (2003e:17023)}


\bibitem{Na-reflect}
Nakajima H.,  Ref\/lection functors for quiver varieties and {W}eyl group
  actions, {\it Math. Ann.} \textbf{327} (2003), 671--721.
  %\MR{MR2023313   (2004k:16036)}


\bibitem{Na-ADHM}
Nakajima H.,  Sheaves on $ALE$ spaces and quiver varieties, {\it Mosc. Math. J.}
  \textbf{7} (2007),  699--722.


\bibitem{Saito}
Saito Y., Crystal bases and quiver varieties, {\it Math. Ann.} \textbf{324}
  (2002),  675--688, \href{http://arxiv.org/abs/math.QA/0111232}{math.QA/0111232}. %\MR{MR1942245 (2004a:17023)}



\end{thebibliography}
\end{document}
